\documentclass{article}
\usepackage[utf8]{inputenc}
\usepackage[T1]{fontenc}
\usepackage[a4paper,total={6in,8in}]{geometry}
\usepackage{amsmath,amssymb,amsfonts}
\usepackage{graphicx}
\usepackage{hyperref}
\usepackage{enumitem}
\usepackage{listings}
\usepackage{xcolor}
\DeclareMathOperator{\Var}{Var}
\DeclareMathOperator{\Cov}{Cov}
\providecommand{\tightlist}{\setlength{\itemsep}{0pt}\setlength{\parskip}{0pt}}
\lstset{basicstyle=\ttfamily\scriptsize,breaklines=true,columns=fullflexible,
  keepspaces=true,showstringspaces=false,frame=single}
\hypersetup{colorlinks=true,urlcolor=blue,linkcolor=blue}
\setlength{\emergencystretch}{2em}

\title{Math 126 Problem 21: Sketched Regression and the Fast JL Transform}
\author{} % Put your name here.
\date{Spring 2027}
\begin{document}
\maketitle

\hypertarget{submission-instructions}{%
\subsection*{Submission Instructions}\label{submission-instructions}}

Please upload your work to Gradescope by \textbf{11:59pm on Monday, May 3, 2027 (Pacific time).}

\begin{itemize}
\item
  You are encouraged to discuss ideas and work with your classmates. However, you \emph{must acknowledge} your collaborators at the top of each solution on which you collaborated with others and you \emph{must write} your solutions independently.
\item
  Your solutions to theory questions must be written legibly, or typeset in LaTeX or markdown.
\item
  I recommend that you write your solutions in a Python notebook using Google Colab.
\item
  You should submit your solutions as a \emph{single PDF} via the assignment on Gradescope.
\end{itemize}

\textbf{Grading:} The point of the problem set is for \emph{you} to learn. To this end, I hope to disincentivize the use of LLMs by \emph{not} grading your work for correctness. Instead, you will grade your own work by comparing it to my solutions. This self-grade is due the Friday \emph{after} the problem set is due, also on Gradescope.

\noindent Original notebook: \href{https://colab.research.google.com/drive/1W2OZ-UWNjreSrctrVbRIce-wItRvGSBZ?usp=sharing}{Google Colab}.

% Source notebook cell 2
\hypertarget{sketched-regression-and-the-fast-jl-transform}{%
\subsection*{Sketched Regression and the Fast JL Transform}\label{sketched-regression-and-the-fast-jl-transform}}

Consider the linear regression problem with \(n \geq d\). Let \(\mathbf{A} \in \mathbb{R}^{n \times d}\) be a feature matrix and \(\mathbf{b} \in \mathbb{R}^n\) be a target vector. The regression problem is to find a minimizing vector \begin{align*}
\mathbf{x}^* = \arg \min_{\mathbf{x} \in \mathbb{R}^d}
\| \mathbf{Ax} - \mathbf{b} \|_2^2.
\end{align*}

You previously showed that the optimal solution is \(\mathbf{x}^* = \left(\mathbf{A}^\top \mathbf{A} \right)^{-1} \mathbf{A}^\top \mathbf{b}\). In this problem, you will compare computing the optimal solution exactly to computing it approximately using the fast Johnson-Lindenstrauss transform. We will use the MNIST dataset to build \(\mathbf{A}\) and \(\mathbf{b}\). The MNIST dataset consists of \(28 \times 28\) pixel handrawn digits of numbers with the corresponding label.

\hypertarget{part-a-standard-solution}{%
\subsection*{Part A: Standard Solution}\label{part-a-standard-solution}}

Load a regression problem of your choice and compute the exact solution \(\mathbf{x}^*\) and the mean squared error \begin{align*}
    \frac1{n} \| \mathbf{Ax} - \mathbf{b} \|_2^2.
\end{align*}

If your code is anything like mine, it will be slow and return a a terrible solution due to \textit{round off error}.

\hypertarget{part-b-fast-jl-transform}{%
\subsection*{Part B: Fast JL Transform}\label{part-b-fast-jl-transform}}

Now implement the fast JL transform as described in class. In particular, compute \(\mathbf{\Pi A} = \mathbf{S H D A}\) one column of \(\mathbf{A}\) at a time. Recall that \(\mathbf{S}\) is a sampling matrix, \(\mathbf{H}\) is a Hadamard, and \(\mathbf{D}\) is a diagonal matrix with a random sign.

When you are done, compute the mean squared error of your solution and comment on how it compares to the ``exact'\,' solution.

\textbf{Hint:} Computing \(\mathbf{H}\) is too expensive so write a function to compute \(\mathbf{HDx}\) using recursion.

\end{document}
